
\documentclass[11pt]{article}
\usepackage[a4paper,margin=0.95in]{geometry}
\usepackage{fontspec}
\setmainfont{Noto Serif}
\usepackage{microtype}
\usepackage{amsmath,amssymb,mathtools}
\usepackage{array,longtable,booktabs}
\usepackage{graphicx}
\usepackage{csquotes}
\usepackage{hyperref}
\hypersetup{colorlinks=false,pdfborder={0 0 0}}
\setlength{\parindent}{1.15em}
\setlength{\parskip}{0.18em}
\allowdisplaybreaks
\sloppy
\newcommand{\scn}[1]{\textsc{#1}}
\newcommand{\W}{\mathrm{W}}
\begin{document}

\begin{center}
{\Large\bfseries BRIEFWECHSEL ZWISCHEN LEJEUNE DIRICHLET UND GAUSS\footnote{Von den Antworten von Gauss auf die im Folgenden mitgetheilten Briefe Dirichlet's standen nur die beiden Briefe zur Verfügung, welche in Gauss Werke Bd. II S. 514--518 abgedruckt sind. Nach einer Mittheilung von Herrn E. Schering haben sich auch im Nachlasse von Gauss andere Antwortschreiben nicht gefunden. F.}\par}
\vspace{0.8em}
\end{center}

\section*{I. Lejeune Dirichlet an Gauss}
\begin{center}\emph{Verehrungswürdiger Herr Hofrath!}\end{center}

Ew. Hochwohlgeboren habe ich die Ehre, meinen ersten mathematischen Versuch als ein Zeichen meiner innigsten Verehrung und meiner tiefsten Bewunderung ganz gehorsamst zu überreichen. Bei der gütigen Empfehlung des Herrn Baron von \scn{Humboldt} glaube ich mir mit der Hoffnung schmeicheln zu dürfen, dass Sie diese erste Arbeit eines jungen Deutschen mit gütiger Nachsicht aufnehmen und beurtheilen werden. Wenn Ew. Hochwohlgeboren dieselbe Ihrer Aufmerksamkeit nicht ganz unwürdig finden sollten, so würde ich es wagen, Sie ganz ergebenst um die Erlaubniss zu ersuchen, Ihnen zuweilen schreiben und Sie um einige Winke zur Leitung meiner ferneren wissenschaftlichen Bestrebungen bitten zu dürfen. Ich würde diese Erlaubniss als das höchste Glück ansehen, indem ich bei der Vorliebe, womit ich das Studium der unbestimmten Analysis betreibe, nichts so sehnlich wünsche, als den Verfasser der unsterblichen \emph{Disquisitiones arithmeticae} an meinen Bemühungen einigen Antheil nehmen zu sehen. Indem ich mich bisher hauptsächlich mit der höheren Arithmetik beschäftigt habe, bin ich ganz meiner Neigung gefolgt, ohne gehörig zu erwägen, wie wenig ich bei meinen beschränkten Anlagen hoffen darf, in diesem schwierigen Theile der Mathematik etwas Erhebliches zu leisten. Obgleich ich mich nun täglich mehr von den Schwierigkeiten überzeuge, womit Untersuchungen dieser Art verbunden sind, so ist mir doch diese Beschäftigung zu sehr zur Gewohnheit und ich möchte fast sagen, zu sehr zur Leidenschaft geworden, als dass ich mich so leicht entschliessen könnte, die einmal begonnenen Untersuchungen aufzugeben.

Gegen das Ende dieses Jahres werde ich nach meinem Vaterlande Preussen zurückkehren, wo ich durch die gütige Verwendung des Herrn Baron von \scn{Humboldt} angestellt zu werden hoffe. Obgleich ich nun die Fürsprache dieses grossen Gelehrten gehörig zu schätzen weiss, so kann ich doch nicht umhin zu gestehen, dass dieselbe mir noch etwas zu wünschen übrig lässt. Es hängt nämlich grossentheils von dem Urtheile, welches man in Berlin über meine Arbeit fällen wird, ab, was für ein Wirkungskreis mir in meinem Vaterlande wird angewiesen werden, und es ist daher bei dem geringen Werthe meines Versuchs von der höchsten Wichtigkeit für mich, dass man denselben mit Nachsicht beurtheile oder demselben doch wenigstens Gerechtigkeit widerfahren lasse. Da aber selbst unter den ausgezeichnetesten Mathematikern -- wie ich mich hier zu überzeugen Gelegenheit gehabt habe -- nur sehr wenige mit der unbestimmten Analysis vertraut sind, so steht zu fürchten, dass meine Abhandlung eine weniger günstige Aufnahme finden werde, als derselben vielleicht zu Theil werden dürfte, wenn dieselbe bei übrigens gleichem inneren Werthe eine Aufgabe aus einem bekanntern Theile der Wissenschaft, z. B. der Astronomie oder Integralrechnung zum Gegenstande hätte. Dieser Umstand wird mich einigermassen entschuldigen, dass ich es wage Ew. Hochwohlgeboren mit der unterthänigsten Bitte zu beschweren, gelegentlich einem der Berliner Gelehrten, mit denen Sie in Correspondenz stehen, Ihr Urtheil über meine Arbeit mitzutheilen. Ich glaube der Hoffnung Raum geben zu dürfen, dass Ew. Hochwohlgeboren diese Gunst einem jungen Deutschen nicht versagen werden, dessen Glück durch Ihre gütige Fürsprache so sehr befördert werden würde.

Da ich noch einige Zeit in Paris zu bleiben gedenke, so bin ich so frei, Ihnen meine Dienste während meines Aufenthaltes in hiesiger Stadt anzubieten, indem ich Ew. Hochwohlgeboren zugleich bitte, wenn Sie mich mit einem Auftrage oder einigen Zeilen beehren wollen, Ihren Brief gefälligst an meinen Vater, Postkommissar in Düren bei Aachen, zu adressiren.

\begin{flushright}
Ich verharre in der tiefsten Verehrung\\[0.4em]
Ew. Hochwohlgeboren\\
ganz gehorsamer\\
\scn{G. Lejeune Dirichlet}
\end{flushright}
Paris d. 28. Mai 1826.

\section*{II. Gauss an Lejeune Dirichlet}
\begin{flushleft}
A Monsieur\\
\hspace*{2em}Monsieur \scn{Lejeune Dirichlet} à Paris.
\end{flushleft}

Schon früher würde ich Ihnen meinen Dank für die mir gütigst übersandte Abhandlung und das grosse Vergnügen, welches Sie mir dadurch gemacht haben, bezeugt haben, wenn ich nicht gewünscht hätte, erst etwas von dem Erfolg dessen zu erfahren, was ich in Beziehung auf Ihre, und ich kann hinzusetzen meine eigenen Wünsche in Berlin zu thun versucht habe. Ich freue mich ungemein jetzt aus einem von dem Secretair der Akademie in Berlin erhaltenen Briefe zu sehen, dass wir hoffen können, dass man Ihnen bald im Vaterlande eine angemessene Fixirung zu verschaffen geneigt sein wird.

Es ist mir eine um so erfreulichere Erscheinung, dass Sie mit grosser Neigung demjenigen Theile der Mathematik anhängen, der von jeher mein Lieblingsstudium gewesen ist, je seltener dieselbe ist. Ich wünsche Ihnen herzlich eine äussere Lage, wo Sie soviel als möglich Herr Ihrer Zeit und der Wahl Ihrer Arbeiten bleiben. Ich selbst wurde gleich nach dem Erscheinen meiner \emph{Disquisitiones} durch andersartige Beschäftigungen, und später durch meine äusseren Verhältnisse sehr gehindert, meiner Neigung in dem Maasse nachzuhängen, wie ich gewünscht hätte. Anstatt eines zweiten Theils jenes Werks, den ich früher beabsichtigte, werde ich mich aller Wahrscheinlichkeit nach darauf beschränken müssen, von Zeit zu Zeit ein Memoire über einen einzelnen Gegenstand zu liefern. Die drei Abhandlungen dieser Art, die bisher im 16. Band der hiesigen Commentationen, und im ersten und vierten der \emph{Commentationes recentiores} erschienen sind, enthalten aber (einen Theil der zweiten abgerechnet) keine von den Gegenständen, die ich schon 1801 zur Fortsetzung im Auge hatte, sondern neue; und so beziehen sich auch meine späteren Arbeiten dieser Art gleichfalls auf einen neuen Gegenstand, namentlich die Theorie der biquadratischen Reste, die ich etwa in drei Abhandlungen zu geben denke; die erste davon wird in kurzem für den sechsten Band der \emph{Comment. rec.} gedruckt werden, und die Hauptmaterialien für das Uebrige sowie für die ähnliche Theorie der cubischen Reste, ist, obgleich noch wenig davon ordentlich zu Papier gebracht ist, im Wesentlichen als abgemacht zu betrachten.

Empfehlen Sie mich gefälligst dem Herrn von \scn{Humboldt}, falls er noch in Paris ist, und entschuldigen mich, dass ich jetzt nicht an ihn selbst schreibe, mit der Besorgniss, dass mein Brief ihn nicht treffen möchte, da er, wie ich höre, Paris zu verlassen die Absicht hatte.

\begin{flushright}
Mit aufrichtiger Hochschätzung\\
Ihr ergebenster\\
\scn{C. F. Gauss}
\end{flushright}
Göttingen, den 13. September 1826.

\section*{III. Lejeune Dirichlet an Gauss}
\begin{center}\emph{Höchstzuverehrender Herr Hofrath!}\end{center}

Schon früher würde ich von der mir von Ew. Hochwohlgeboren gütigst ertheilten Erlaubniss, Ihnen zuweilen schreiben zu dürfen, Gebrauch gemacht haben, um Ihnen für das Wohlwollen, womit Sie mich aufgenommen haben, zu danken, wenn ich nicht gewünscht hätte, Ihnen zugleich eine Arbeit zu übersenden, mit der ich mich bald nach meiner Ankunft in Breslau zu beschäftigen angefangen hatte. Diese Arbeit ist durch die in Ihren Anzeigen enthaltene Ankündigung Ihrer Abhandlung über die biquadratischen Reste veranlasst worden. Die Eleganz der in dieser Ankündigung mitgetheilten Kriterien zur Entscheidung der Frage, ob $\pm2$ biquad. Rest oder Nichtrest einer Primzahl $=8n+1$ sey, deren zweites mir ganz besonders merkwürdig schien, erregte bei mir den Wunsch, meinerseits Beweise für dieselben zu finden. Sobald ich mich hier mit meinen Berufsgeschäften etwas vertraut gemacht hatte, fing ich an, mich mit diesem Gegenstande ernstlich zu beschäftigen; meine Bemühungen blieben eine Zeit lang fruchtlos, bis es mir gelang, Ihr zweites Kriterium, dessen Beweis nach dem von Ihnen gewählten Gange eine Menge Hülfsuntersuchungen zu erfordern scheint, aus dem ersten abzuleiten. Allein nach diesem ersten glücklichen Erfolge gerieth meine Arbeit wieder sehr ins Stocken und ich konnte lange Zeit kein zur Begründung des ersten Kriteriums geeignetes Mittel ausfindig machen. Endlich kam ich gegen den Anfang des Winters auf den in der beiliegenden Abhandlung enthaltenen Beweis, der so einfach ist, dass man kaum begreift, wie sich derselbe einem nicht gleich darbietet, sobald man nur mit dem zu beweisenden Satze bekannt ist. Die Fortsetzung meiner Untersuchungen, von denen meine Abhandlung nur einen Theil enthält, hat mich auf eine Menge von Resultaten geführt, von denen man im voraus gewiss nicht vermuthet hätte, dass sie mit diesem Gegenstande in irgend einer Beziehung ständen, und es haben sich mir bei dieser Arbeit mehrere merkwürdige Beispiele von dem oft wunderbaren Zusammenhange arithmetischer Wahrheiten dargeboten, in welchem Sie den Hauptgrund des Reizes finden, welchen uns die Untersuchungen der unbestimmten Analysis gewähren.

Da ich voraussetzen konnte, dass die Leser meiner Abhandlung schon einigermassen mit der unbestimmten Analysis vertraut seyn würden, so habe ich mich bei der Abfassung der Preliminarien sehr kurz gefasst. Dieser Umstand hat eine der Allgemeinheit der Untersuchung nachtheilige Beschränkung veranlasst. Ich habe mich nämlich durch Analogie zu dem Irrthume verleiten lassen, als sey es in der Theorie der biquad. Reste hinlänglich, das Verhalten irgend einer Primzahl zu einer andern Primzahl der Form $4n+1$ auszumitteln, da doch zur vollständigen Beantwortung aller hierher gehörigen Fragen erfordert wird, dass man auch das Verhalten von zusammengesetzten Zahlen, wenigstens von solchen die Produkte zweier Primzahlen sind, zu Primzahlen der Form $4n+1$ bestimme. So kann z. B. die Frage, ob $21$ biq. Rest oder Nichtrest einer Primzahl $84n+1$, $25$, $37$ sey, sehr leicht auch vermittelst der in der Abhandlung aufgestellten Sätze beantwortet werden; anders aber verhält es sich mit den Primzahlen $p$ der Form $84n+5$, $17$, $41$, welche neue Kriterien erfordern. Für den eben erwähnten speciellen Fall finde ich folgendes Kriterium.

\begin{quote}
Setzt man $5p=\varphi^2+\psi^2$ (wo ich annehme, dass $\psi$ durch $6$ theilbar sey), so wird $21$ biq. Rest von $p=84n+5$, $17$, $41$ seyn, wenn eine der Zahlen $\varphi$, $\psi$ durch $7$ theilbar; im entgegengesetzten Falle ist $21$ biq. Nichtrest von $p$.
\end{quote}

Ich bin jetzt eifrig damit beschäftigt, die durch den oben angegebenen Irrthum in meiner Arbeit hervorgebrachte Lücke auszufüllen.

Sie hatten während meines Aufenthaltes in Göttingen die Güte, mir einige der merkwürdigen Resultate mitzutheilen, welche die Frucht der von Ihnen neuerdings angestellten geometrischen Forschungen sind. Der schon damals entstandene Wunsch, Ihre Untersuchungen kennen zu lernen, ist recht lebhaft bei mir geworden, seitdem ich durch den in den Anzeigen enthaltenen Auszug aus Ihrer Abhandlung von den in der Unterhaltung nur leicht angedeuteten Resultaten eine deutlichere und vollständigere Idee bekommen habe. Es ist mir bei dieser Gelegenheit wieder recht fühlbar geworden, wie sehr es zu wünschen wäre, dass die gelehrten Gesellschaften bei der Herausgabe ihrer Commentarien dem von der Londoner Societät gegebenen Beispiele folgen möchten, die jedes halbe Jahr einen Band erscheinen lässt. Sogar Ihre Abhandlung über die biquadratischen Reste ist mir noch immer nicht zugekommen, obgleich ich die schleunigste Besorgung derselben dem Buchhändler zu wiederholten Malen anempfohlen habe und der Band Ihrer Commentarien, welcher dieselbe enthalten soll, schon im vorigen Katalog angekündigt ist.

Dürfte ich Ew. Hochwohlgeboren bitten, die Versicherung der innigen Verehrung und Dankbarkeit zu genehmigen, womit ich verharre

\begin{flushright}
Ew. Hochwohlgeboren\\
gehorsamer\\
\scn{Lejeune Dirichlet}
\end{flushright}
Breslau d. 8. April 1828.

\section*{IV. Gauss an Lejeune Dirichlet}

Für Ihr gütiges Schreiben, und die gefällige Uebersendung Ihrer beiden Abhandlungen statte ich Ihnen, mein hochgeschätzter Freund, meinen verbindlichsten Dank ab. Ich sehe mit Vergnügen das steigende Interesse, welches man gegenwärtig an den Untersuchungen der Höhern Arithmetik zu nehmen anfängt. Die glückliche Art, wie Sie das zweite auf die biquadratische Residualität der Zahl $2$ [bezügliche Kriterium] aus dem ersten ableiten, hat mir sehr wohl gefallen.

Vermuthlich hat jetzt der 6. Band unserer Commentationen seinen Weg nach Breslau gefunden, und meine \emph{Commentatio prima} über die biquadratischen Reste wird Ihnen also wohl gegenwärtig bekannt sein: wenn sich eine Gelegenheit darbieten sollte, würde ich auch mit Vergnügen Ihnen einen besonderen Abdruck derselben übersenden. Ich hätte unter mehreren Beweisarten für das darin vorkommende Theorem wählen können; es wird Ihnen aber nicht entgehen, warum ich den daselbst ausgeführten hier vorgezogen habe, hauptsächlich nemlich, weil die Classification von $2$ bei denjenigen Moduln, für welche es quadratischer Nichtrest ist (unter $B$ oder $D$) als ein wesentlicher integrirender Theil des Theorems betrachtet werden muss, auf welchen die meisten andern Beweisarten nicht anwendbar scheinen.

Die ganze Untersuchung, deren Stoff ich schon seit 23 Jahren vollständig besitze, die Beweise der Haupttheoreme aber (zu welchen das in der ersten Commentation noch nicht zu rechnen ist) seit etwa 14 Jahren -- (obwohl ich wünsche und hoffe, an letztern, den Beweisen, noch einiges vereinfachen zu können) -- habe ich auf ungefähr drei Abhandlungen berechnet. Mit der Abfassung der zweiten habe ich bereits jetzt einen Anfang gemacht, und hoffe, sie in nicht langer Zeit zu vollenden, falls nicht die neuerdings mir wieder aufgetragenen Messungsgeschäfte dabei noch einige Verzögerung verursachen.

Das Schlusstheorem $b\equiv \frac12 rr'\pmod p$ hatte ich schon vor drei Jahren in den hiesigen gel. Anzeigen mit bekannt, und auf den merkwürdigen dabei noch zu lösenden Knoten aufmerksam gemacht; ich habe aber bisher nicht gehört, dass jemand einen Versuch dazu gemacht hätte. Vor einigen Tagen ist es mir nun mit der einen Hälfte wirklich gelungen, und dieser Fund hat mir um so mehr Vergnügen gemacht, da er sich garnicht auf Induction gründet -- denn ich gestehe, dass ich gerade diesen Zusammenhang nicht erwartet hätte sondern a priori auf die Combination anderweitiger sehr verschlungener und interessanter, schon 28 Jahr alter, aber noch garnicht bekannt gemachter Untersuchungen, wovon eine leise Andeutung in der Schlussanmerkung der \emph{Disquis. Arithm.} S. 668\footnote{Gauss Werke Bd. I. S. 466.} gegeben ist.

Es ist dies nemlich ein ausreichendes Criterium für den Fall, wo $p$ von der Form $8n+5$ ist.

Es sei die Anzahl der Klassen, welche die binären Formen in jeder der beiden Gattungen für den Determinant $-p$ bilden gleich $k$. Der Anfang einer von mir bis zu dem Determinant $-3000$ construirten Tafel steht \emph{Disquis. Arithm.} p. 520\footnote{art. 303.}. Auch ist noch zu bemerken, dass für ein $p$ von der angenommenen Form, allemal $k=2m+1$ wird, wenn $m$ die Anzahl der Zerlegungen von $p$ in drei positive Quadrate bedeutet (ich sage positiver, um 0 auszuschliessen), wie \scn{Legendre} durch Induction gefunden, und in den \emph{Disquisitiones Arithm.} zuerst aus der Theorie der ternären Formen bewiesen ist. Man hat z. B.

\begin{center}\scriptsize
\begin{tabular}{c|cccccccccc}
$p$&5&13&29&37&53&61&101&109&149&157\\
$k$&1&1&3&1&3&3&7&3&7&3\\
$m$&0&0&1&0&1&1&3&1&3&1\\
Zerlegungen&&&$4+9+16$&&$1+16+36$&$9+16+36$&\begin{tabular}{c}$1+36+64$\\$4+16+81$\\$16+36+49$\end{tabular}&$9+36+64$&\begin{tabular}{c}$1+4+144$\\$4+64+81$\\$36+64+49$\end{tabular}&$4+9+144$
\end{tabular}\end{center}

Dies vorausgesetzt, ist allemal derjenige Werth von $b$, welcher $\equiv\frac12 rr'\pmod p$ ist,
\[
 b\equiv 2k+a-1\equiv 4m+a+1 \pmod 8,
\]
wodurch das Zeichen von $b$ vollkommen bestimmt ist. Sehen Sie hier 22 Beispiele, indem ich die Ausdehnung der am Schluss der Abhandlung gegebenen Tafel verdopple:

\begin{center}\small
\begin{tabular}{rrrr|rrrr}
\toprule
$p$&$k$&$a$&$b$&$p$&$k$&$a$&$b$\\
\midrule
5&1&$+1$&$+2$&181&5&$+9$&$+10$\\
13&1&$-3$&$-2$&197&5&$+1$&$-14$\\
29&3&$+5$&$+2$&229&5&$-15$&$+2$\\
37&1&$+1$&$-6$&269&11&$+13$&$+10$\\
53&3&$-7$&$-2$&277&3&$+9$&$+14$\\
61&3&$+5$&$-6$&293&9&$+17$&$+2$\\
101&7&$+1$&$-10$&317&5&$-11$&$+14$\\
109&3&$-3$&$+10$&349&7&$+5$&$+18$\\
149&7&$-7$&$-10$&373&5&$-7$&$+18$\\
157&3&$-11$&$-6$&389&11&$+17$&$-10$\\
173&7&$+13$&$+2$&397&3&$-19$&$-6$\\
\bottomrule
\end{tabular}\end{center}

Man kann die Vorschrift also auch so ausdrücken, (immer voraussetzend $p\equiv5\pmod 8$). Es ist
\[
 b\equiv a+1 \pmod 8,\quad \text{wenn }m\text{ gerade},
\]
\[
 b\equiv a+5 \pmod 8,\quad \text{wenn }m\text{ ungerade}.
\]
Ich wage noch keine Vermuthung, ob ein noch einfacheres Criterium möglich ist, woran man den Fall des geraden $m$ von dem des ungeraden im Voraus unterscheiden könnte, d. i. ohne den Werth von $m$ selbst zu kennen, da, wie ich schon oben bemerkt habe, dies Rapprochement noch ganz neu ist. Für den Fall $p\equiv1\pmod 8$, bleibt zwar obige Congruenz $b\equiv2k+a-1\pmod 8$ richtig, entscheidet aber nicht mehr über das Zeichen von $b$, da sie dem positiven und negativen Werthe von $b$ zugleich genug thut. Es ist hier nemlich $k$ immer gerade, $=2m$ (wenn die Bedeutung von $m$ ebenso ausgesprochen wird wie oben) oder $=2m+2$, wenn man unter $m$ die Anzahl der Zerlegungen von $p$ in 3 positive ungleiche Quadrate versteht, und $b\equiv0\pmod4$, oder $b\equiv -b\pmod8$. Ich vermuthe, dass der Fall $p\equiv1\pmod8$ oder $b\equiv0\pmod4$ \emph{altioris indaginis} ist und vielleicht wieder
\[
 b\equiv4\pmod 8\quad\text{leichter als}\quad b\equiv0\pmod8,
\]
\[
 b\equiv8\pmod {16}\quad\text{leichter als}\quad b\equiv0\pmod{16}\quad \text{u. s. w.}
\]

\begin{flushright}
Mit ausgezeichneter Hochachtung beharre ich\\
Ihr freundschaftlich ergebenster\\
\scn{C. F. Gauss}
\end{flushright}
Göttingen, den 30. Mai 1828.

\section*{V. Lejeune Dirichlet an Gauss}

Empfangen Sie, hochgeehrter Gönner, meinen innigsten Dank für die mir gütigst ertheilte Erlaubniss, einige Zeit in Ihrer Nähe zubringen zu dürfen. Leider bin ich nicht im Stande, in diesen Ferien davon Gebrauch zu machen und einen lange gehegten Wunsch auszuführen. Ich habe vor wenigen Tagen einen ziemlich heftigen Ruhranfall zu bestehen gehabt, und jetzt wo ich mich wieder herzustellen anfange, liegen die Meinigen an den Masern danieder. Ich darf Ihnen wohl nicht erst sagen, wie sehr ich es bedauere, meinen Besuch in Göttingen abermals hinausgeschoben zu sehen. Seit Jahren mit dem Studium Ihrer Werke beschäftigt, könnte für mich nichts genussbringender und fördernder seyn, als Ihrer Güte Aufschlüsse über die Entstehung und Ausbildung der herrlichen Methoden zu verdanken, welche in Ihren Schriften in bewundernswürdiger Vollendung erscheinen. Ja, ich will es Ihnen nur gestehen, ich gehe in meinen Hoffnungen so weit, mir von Ihrem Wohlwollen mündliche Mittheilung eines Theiles Ihrer theoretischen Untersuchungen über den Magnetismus und Elektromagnetismus zu versprechen, und gelobe im voraus feierlich, wenn Sie mich einer solchen Gunst würdigen sollten, mir Herrn \scn{Krone}'s Diskretion zum Vorbilde zu nehmen, welcher im beneidenswerthen Besitz dieses Schatzes nichts hat verrathen wollen, als dass sich diese Untersuchungen an die schöne in Ihrer Abhandlung über die Attraktion der Ellipsoide entwickelte Methode innig anschliessen. Haben Sie wohl durch die \emph{Comptes rendus} der Pariser Akademie von der Diskussion Kenntniss genommen, welche sich neuerdings über das eben erwähnte Problem erhoben hat. \scn{Poisson} scheint seine Auflösung für die einzige direkte zu halten, während bei allen andern im Falle eines äussern Punktes die Schwierigkeit nur eludirt werde. Ich muss aufrichtig bekennen, dass dergleichen Aeusserungen für mich keinen rechten Sinn haben. Ein doppeltes Integral durch Einführung neuer Variabeln, für welche sich die eine Integration ausführen lässt, auf ein einfaches zurückführen, scheint mir nicht direkter als irgend ein anderes Mittel z. B. die Differentiation unter dem Integralzeichen. Nach meiner Ansicht kommt Alles darauf an, auf welchem Wege sich die Zurückführung am schnellsten und übersichtlichsten bewerkstelligen lässt, und in dieser Beziehung hat wohl Ihre Behandlung, welche die Attraktion auf innere und äussere Punkte auf einen Schlag giebt, entschieden den Vorzug vor allen bekannten. Wollte man dergleichen Distinktionen gelten lassen, so müsste man, um consequent zu seyn, bei fast allen bestimmten Integralen zugeben, dass bei ihrer Auffindung die Schwierigkeit nur eludirt sey, womit jedoch, wie mir scheint, eigentlich gar nichts gesagt wäre.

Erlauben Sie, hochgeehrter Gönner, dass ich Ihnen beiliegend einen kleinen im letzten Hefte des Crelleschen Journals erschienenen Aufsatz überreiche. Ich bin auf die darin angedeutete, mir sehr fruchtbar scheinende Methode durch die Bemühung geführt worden, den Satz, dass jede arithmetische Progression, deren erstes Glied und Differenz keinen gemeinschaftlichen Factor haben, unendlich viel Primzahlen enthält, zu beweisen. Ich möchte fast vermuthen, dass meine Methode mit den in der Schlussbemerkung der \emph{disq. arith.} angedeuteten Untersuchungen einige Analogie hat, besonders deshalb, weil Sie von Ihren Untersuchungen sagen, dass sie über mehrere Theile der Analysis Licht verbreiten, und meine Methode in so innigem Zusammenhange mit den merkwürdigen trigonometrischen Reihen steht, welche discontinuirliche Funktionen darstellen und deren Natur zur Zeit des Erscheinens der \emph{disq. arith.} noch ganz unaufgeklärt war. Sie werden sich vielleicht noch erinnern, dass Sie mir vor mehr als 10 Jahren, als ich noch in Breslau war, ein Kriterium zur Entscheidung der am Ende Ihrer ersten Abhandlung über biq. Reste aufgeworfenen Frage (für den Fall, wo $p=8n+5$) mit der Bemerkung mitgetheilt haben, dass Sie dasselbe aus den am Ende der \emph{disq. arith.} angekündigten Untersuchungen abgeleitet hätten. Dieses Kriterium folgt nun ganz leicht aus der Combination des Satzes, dass für die Determinante $-p$ die Anzahl der Formen in jedem Genus dem Ueberschuss der Anzahl der quad. Reste $<\frac p4$ über die Anzahl der quad. Nichtreste $<\frac p4$ gleich ist, mit der Theorie der biquad. Reste, diese Theorie selbst nur in dem Umfange genommen, worin ich dieselbe im Crelleschen Journal gegeben habe. Darf ich wohl bei einer ausführlichen Ausarbeitung meiner Untersuchungen über den Gebrauch der Reihen in der höheren Arithmetik, womit ich mich nächstens zu beschäftigen gedenke, dieses Kriterium als mir von Ihnen im Jahr 1828 mitgetheilt erwähnen?

Erlauben Sie mir die Bitte, beiliegenden Brief \scn{Weber} einhändigen zu wollen und genehmigen Sie gefälligst die Versicherung der tiefsten Bewunderung und innigsten Verehrung.

\begin{flushright}
Ihres dankbar ergebenen\\
\scn{Dirichlet}
\end{flushright}
Berlin d. 9. Sept. 1838.

\section*{VI. Lejeune Dirichlet an Gauss}
\begin{center}\emph{Verehrtester Gönner!}\end{center}

Die bei mir durch die Kränklichkeit, wovon meine Frau in der letzten Zeit heimgesucht worden ist, verursachte Verstimmung wird mich hoffentlich entschuldigen, dass ich so lange gezögert habe, Ihnen zu schreiben, wie sehr ich auch wünschte, Ihnen für die überaus gütige und wohlwollende Aufnahme, deren Sie mich gewürdigt haben, meinen innigsten Dank abzustatten. Eine andere Veranlassung meinen Vorsatz zu verschieben, lag in dem Umstande, dass mehrere der Bücher, welche ich zu Rathe ziehen wollte, um Ihnen die gewünschte Auskunft über das zu geben, was bisher über die kürzeste Fläche geschrieben worden, gerade nicht auf der Bibliothek waren, und ich daher die Zurücklieferung derselben erst abwarten musste.

Ueber den genannten Gegenstand habe ich nun zwar gar nichts gefunden, wohl aber Untersuchungen über andere Probleme, die wenn auch dem Objekte nach verschieden, doch im Wesen mit der Bestimmung der kürzesten Fläche zusammenfallen. Diese Probleme betreffen die Theorie der Wärme für den Fall permanenter Temperaturen, auf den sich je nachdem man 2 oder 3 Dimensionen in Betracht zieht, die Gleichung
\[
(1)\quad \frac{\partial^2u}{\partial x^2}+\frac{\partial^2u}{\partial y^2}=0,
\qquad\text{oder}\qquad
(2)\quad \frac{\partial^2u}{\partial x^2}+\frac{\partial^2u}{\partial y^2}+\frac{\partial^2u}{\partial z^2}=0,
\]
bezieht. Projicirt sich die gegebene von einer Ebene sich wenig entfernende Curve, durch welche die kürzeste Fläche gelegt werden soll, als Rechteck oder Kreis, so ist die Auflösung der Aufgabe leicht aus den von \scn{Fourier} in seiner \emph{Théorie de la Chaleur} gegebenen Methoden abzuleiten. Auch der weit schwierigere Fall, wo die Curve sich als Ellipse projicirt, scheint implicite in einer schönen Abhandlung von \scn{Lamé} (Journal de Liouville, Mai 1839) enthalten, indem dort eine Funktion $u$ bestimmt wird, die innerhalb eines durch die Gleichung
\[
 \left(\frac{x}{\alpha}\right)^2+\left(\frac{y}{\beta}\right)^2+\left(\frac{z}{\gamma}\right)^2=1
\]
gegebenen Ellipsoides der Differentialgleichung (2) genügt und sich auf der Oberfläche auf eine beliebige Funktion der beiden jeden Punkt derselben bestimmenden Coordinaten reducirt. Setzt man $\gamma=\infty$, und lässt die gegebene beliebige Funktion von $z$ unabhängig werden, so erhält man die Gleichung der kürzesten Fläche, deren Grenzcurve sich als Ellipse projicirt.

Ich habe mich seit meiner Rückkunft mit den quadratischen Formen für complexe Zahlen etwas beschäftigt und mich überzeugt, dass fast alle Resultate und Methoden der 5ten Sekt. der \emph{disq. arith.} \emph{mutatis mutandis} auf diesen Fall anwendbar bleiben. Auch lässt sich in ähnlicher Weise wie für reelle Zahlen die Anzahl der Formen bestimmen, wobei sich das merkwürdige Resultat ergiebt, dass diese Anzahl, wie sie für reelle positive Determinanten mit der Kreistheilung zusammenhängt, hier in analoger Beziehung zur Theilung der Lemniscate steht. Diese Untersuchungen haben mich auf einen Satz geführt, welcher für die Theorie der unbestimmten Gleichungen höherer Grade nicht unwichtig scheint, und eine merkwürdige Verallgemeinerung des Theorems darbietet, nach welchem die Gleichung $t^2-Du^2=1$ immer auflösbar ist.

Sind $a,b,\ldots,g,h$ ganze Zahlen und hat die Gleichung
\[
 s^n+as^{n-1}+bs^{n-2}+\cdots+gs+h=0
\]
keinen rationalen Faktor und unter ihren Wurzeln $\alpha,\beta,\gamma,\ldots$ wenigstens eine reelle, so lassen sich immer und auf unendlich verschiedene Art ganze Zahlen
\[
 \lambda,\mu,\nu,\ldots,\varrho\qquad(\text{ohne dass }\lambda=1,\ \mu=0,\ \nu=0,\ldots)
\]
von solcher Beschaffenheit finden, dass, wenn man zur Abkürzung setzt
\[
 \lambda+\mu\alpha+\nu\alpha^2+\cdots+\varrho\alpha^{n-1}=F(\alpha),
\]
das Produkt
\[
 F(\alpha)F(\beta)F(\gamma)\cdots=1
\]
ist. Der Beweis dieses Satzes ist im höchsten Grade einfach und mit Hülfe geometrischer Betrachtungen auf zwei oder drei Seiten abzumachen.

Ich kann diesen Brief nicht schliessen, ohne den schon in Göttingen Ihnen vorgetragenen Wunsch zu wiederholen, dass es Ihnen gefallen möge, mich durch einen besondern Abdruck Ihrer schönen Untersuchungen über die allgemeine Theorie der Attraktionen zu beglücken, welche kennen zu lernen ich mich nur höchst ungern bis zur Vollendung des ganzen nächsten Heftes der \emph{Resultate} etc. gedulden würde. In Erwartung einer gütigen Gewährung meiner Bitte verharre ich

\begin{flushright}
Ihr dankbar ergebenster\\
\scn{Lejeune Dirichlet}
\end{flushright}
Berlin, d. 3 Jan. 1840.

Darf ich Sie wohl bitten, inliegenden Brief \scn{Weber} einzuhändigen oder falls er in Leipzig seyn sollte, dorthin auf die Post geben zu wollen.

\section*{VII. Lejeune Dirichlet an Gauss}
\begin{center}\emph{Hochgeehrter Gönner!}\end{center}

Wenn ich Ihnen das Resultat der in Ihrem Auftrage angestellten Erkundigungen wegen der russischen Bücher so spät mittheile, so wird diese Verspätung, wie ich hoffe, in dem guten Erfolge meiner Bemühungen einige Entschuldigung bei Ihnen finden. Wie Sie selbst hat H. Geheimrath von \scn{Varnhagen} grosse Schwierigkeiten gefunden, sich russische Bücher zu verschaffen. Die früher mit grosser Mühe zu diesem Zweck von ihm angeknüpften Verbindungen bestehen jetzt grossentheils nicht mehr; da er schon seit Jahren seine russischen Studien sehr vernachlässigt und es ihm folglich an Gelegenheit gefehlt hat, die früheren Verbindungen zu cultiviren. In jener früheren Zeit aber, als er das Russische mit grossem Eifer trieb, hat er einen nicht unansehnlichen Vorrath von Büchern erworben, die er Ihnen zu beliebiger Benutzung mit grosser Bereitwilligkeit zur Disposition stellt. Auf meinen Wunsch hat er ein Verzeichniss der wichtigsten dieser Bücher entworfen, welches ich beilege. Was Sie von diesen Werken zu lesen wünschen, bitte ich mir gefälligst zu nennen. Ich werde dann das Verlangte mir sogleich von H. von V. ausbitten und Ihnen ohne Zeitverlust übersenden.

Ich bin gerade zu rechter Zeit von meiner Ferienreise zurückgekehrt, um den armen \scn{Eisenstein} wenigstens noch einmal zu sehen. Als ich ihn den Tag nach meiner Ankunft besuchte, fand ich ihn zum Erschrecken verändert und so ganz mit seinen Leiden beschäftigt, dass selbst der Versuch, ihm von Ihnen zu erzählen und so ihn einen Augenblick zu zerstreuen, ohne Erfolg blieb. Bei so grenzenlosen Leiden musste mir der Tod, der noch in derselben Nacht eintrat, als eine Wohlthat für unsern armen Freund erscheinen.

Wie ich von \scn{Weber} höre, wünschen Sie Auskunft über den Ursprung der mathematischen Bibliothek, welche hier im März verkauft werden soll, wenn sie nicht früher -- wozu einige Wahrscheinlichkeit vorhanden zu seyn scheint -- im Ganzen acquirirt wird. Diese Bibliothek ist früher im Besitz eines Geheimraths \scn{Heiligenstein} gewesen, der sich zu Anfang des Jahrhunderts in Paris aufhielt, und dann später in den seines Sohnes übergegangen, der ein Astronom gewesen seyn soll. Mit dieser Bibliothek hat man zu gleichzeitiger Verauktionirung einige andere Sachen vereinigt, zu denen z. B. die beiden Mst von Ihnen -- wovon das eine das der \emph{Determinatio attract.}, auf den Inhalt des anderen besinne ich mich nicht gleich -- gehören, die aus \scn{Dirksen}'s Nachlass stammen.

Sie schienen neulich zu wünschen, dass mein Beweis für die Convergenz der trigonometrischen Reihen auf den Fall ausgedehnt würde, wo die Max. u. Min. der zu entwickelnden Funktion nicht mehr zählbar sind, und äusserten zugleich die Vermuthung, dass der Nachweis des Theorems in dieser erweiterten Voraussetzung wahrscheinlich keine wesentlichen Schwierigkeiten darbieten würde. Bei näherer Erwägung der Sache finde ich Ihre Vermuthung vollständig bestätigt, wenn man anders von gewissen ganz singulären Fällen absehen will. Die ganze Grundlage des Beweises bilden die Sätze, nach welchen
\[
 \lim \int_0^c f(\beta)\frac{\sin k\beta}{\sin\beta}\,d\beta=\frac{\pi}{2}f(0),
 \qquad\text{und}\qquad
 \lim \int_b^c f(\beta)\frac{\sin k\beta}{\sin\beta}\,d\beta=0,
\]
wo $0<b<c\leq\frac\pi2$ angenommen ist und sich das Zeichen $\lim$ auf das Unendlichwerden der positiven Grösse $k$ bezieht.

Die Erweiterung dieser Sätze auf den eben erwähnten Fall, wo $f(\beta)$ eine unendliche Anzahl von Max. u. Min. darbietet, ist ganz leicht wenigstens für den zweiten Satz und ebenso für den ersten, wenn nur nicht unendlich viele Max. u. Min. in unmittelbarer Nähe von $\beta=0$ liegen. Man darf nämlich aus den Integralen nur die Theile ausscheiden, innerhalb welcher eine solche Anhäufung von Max. u. Min. Statt findet und die man, wenn man die Grenzen dieser Theilintegrale einander hinlänglich nähert, so einrichten kann, dass dieselben beliebig klein bleiben, wie gross auch $k$ werde. Aus dieser einfachen Bemerkung und nach der Art und Weise, wie aus obigen Sätzen die Convergenz geschlossen wird, folgt schon, dass die trigonometrische Reihe zur Darstellung einer beliebigen Funktion $\varphi(x)$ selbst dann, wenn die Max. u. Min. von $\varphi(x)$ nicht mehr zählbar sind, immer convergirt, wofern man $x$ nur keinen Werth beilegt, in dessen unmittelbarer Nachbarschaft eine unendliche Anhäufung von Max. u. Min. von $\varphi(x)$ Statt findet. Um die Convergenz für die eben ausgenommenen Werthe nachzuweisen, muss gezeigt werden, dass der erste der obigen Sätze seine Gültigkeit behält, wenn $f(\beta)$ von $\beta=0$ bis $\beta=\delta$, wo $\delta$ beliebig klein ist, unendlich oft vom Wachsen zum Abnehmen und umgekehrt übergeht. Ist alsdann $f(\beta)$ so beschaffen, dass die Differenz $f(\beta)-f(0)$ sich als ein Produkt darstellen lässt, dessen erster Faktor $\varphi(\beta)$ für ein unendlich kleines $\beta$ endlich bleibt, während der zweite, den ich positiv voraussetze und $\psi(\beta)$ nennen will, die Eigenschaft besitzt, dass das Integral
\[
 \int_0^\delta \psi(\beta)\frac{d\beta}{\beta},
\]
wie es z. B. für jede positive Potenz $\beta^p$ der Fall ist, für ein unendlich kleines $\delta$ selbst unendlich klein wird, so darf man
\[
 \int_0^\delta f(\beta)\frac{\sin k\beta}{\sin\beta}\,d\beta
\]
nur in die beiden Bestandtheile
\[
 f(0)\int_0^\delta \frac{\sin k\beta}{\sin\beta}\,d\beta
 +\int_0^\delta \varphi(\beta)\sin k\beta\frac{\psi(\beta)}{\sin\beta}\,d\beta
\]
zerlegen, von denen der erste für ein unveränderliches $\delta$ durch das Wachsen von $k$ in $\frac\pi2 f(0)$ übergeht, während der zweite für ein gehörig klein gewähltes $\delta$ immer kleiner bleibt als eine beliebig kleine Grösse. Hieraus und da
\[
 \int_\delta^c f(\beta)\frac{\sin k\beta}{\sin\beta}\,d\beta
\]
die Null zur Grenze hat, folgt dann, wenigstens in der vorhin gemachten Voraussetzung über $f(\beta)-f(0)$, die Richtigkeit des ersten Satzes.

Verzeihen Sie, verehrter Gönner, diese lange fast unbriefliche mathematische Exposition, die ich nicht gewagt haben würde, wenn ich eben als ich sie begann die grosse Ausdehnung derselben vorhergesehen hätte.

Genehmigen Sie, hochgeehrter Gönner, den Ausdruck der tiefen Bewunderung und innigen Verehrung.

\begin{flushright}
Ihres dankbar\\
ergebenen\\
\scn{Dirichlet}
\end{flushright}
Berlin 20. Feb. 53.

\end{document}
